\section{Pauli webs of the declared detectors}
\label{app:detector-webs}
We draw all the Pauli semiwebs that correspond to checks, one for each declared detector, in
the full $d=3$ and $d=5$ cultivation circuits of Gidney et al.~\cite{msc}, as provided by
SOFT~\cite{soft}, in the accompanying \texttt{detector\_webs.pdf}.

Some of the declared detectors are not Pauli webs of the $T$ circuit,
Table~\ref{tab:detector-webs}. Their closed semiwebs have defects on $T$ spiders,
Fig.~\ref{fig:detector-web}, so they depend on the phases of individual $T$ gates. An $X$ or $Y$
fault next to a $T$ gate turns it into $T^\dagger$ and can make them random. The other
detectors respond to every Pauli fault in the same deterministic way in the $T$ circuit and the
$S$ proxy. In the proxy every declared detector is a Pauli web, so it flips deterministically
and the proxy always rejects a pattern that flips one of them, while the $T$ circuit may accept
the same pattern. The detectors in Table~\ref{tab:detector-webs} are therefore the only ones that
separate the two, perhaps explaining the lower fault distances.

\begin{table}[h]
\centering
\caption{Declared detectors of the full circuits, numbered from 0 in source order.
Pauli-web detectors hold whatever the phase of each $T$ gate. The others depend on those
phases.}
\label{tab:detector-webs}
\begin{tabular}{lrrrl}
\toprule
 & Declared & Pauli webs & Not Pauli webs & Detectors that are not Pauli webs\\
\midrule
$d=3$ & 20 & 16 & 4 & 7, 14, 16, 18\\
$d=5$ & 107 & 93 & 14 & 7, 19, 20, 23, 69, 89, 91, 93, 95, 97, 99, 101, 103, 105\\
\bottomrule
\end{tabular}
\end{table}

\begin{figure}[h]
\centering
\semiwebinput[.98\linewidth]{figures/detector_webs/d3_det014_check.tikz}
\caption{Closed semiweb for the $d=3$ detector $r_7\oplus r_{15}=0$, the central
readout and the stabilizer $X_0X_3X_7X_{10}$, from the double check to the end of the circuit.
Red, green and blue overlays are $X$, $Z$ and $Y$ labels, and wires cut at the left edge
continue into the injection. The 16 violet stars are defects of $-\pi/2$ at the entry and
readout $T^\dagger$ gates and $+\pi/2$ at the exit and readout $T$ gates. With every $T$
replaced by $S$ the same labels form a Pauli web.}
\label{fig:detector-web}
\end{figure}
